% page 492 du fac-simile (image p-491 du scan)
% bloc 157.5 x 250.6 mm ; marge gauche 8.0 mm
\pgc{-12.700mm}{0.000mm}{
\l{\cel{3}484}
\l{}
\l{\sou{razar}. (\sou{trans}.) Reze-tondar la barbo-pili. - DEFIS.}
\l{}
\l{re. (\sou{muziko}). Noto duesma di la gamo naturala di \textquotedbl{}do\textquotedbl{}. -}
\l{}
\l{\sou{reaktar}. (\sou{netrans., ad, kontre}). I. Facar sur korpo efiko}
\l{\cel{3}genitata da la efiko quan lu recevas de ica. - II. (ke-}
\l{\cel{3}\sou{mio}). (\sou{aludante korpo}). Manifestar sua karakteri distin-}
\l{\cel{3}ganta per la efekto di la efiko quan altra korpo facis a}
\l{\cel{3}lu. - III. Esforcar por rezistar. - DEFIRS.}
\l{}
\l{\sou{reala}. I. Di qua la existo esas fakta, e ne nur semblanta o}
\l{\cel{3}posibla. - II. (\sou{metaf.}) Dicesas pri\cel{2}grandoro, antonime ye}
\l{\cel{3}\sou{imaginara}.) III. (\sou{optiko})\cel{1}\sou{Foko reala}\cel{1}: ube la lumo-radii}
\l{\cel{3}reflektita interrenkontras (antonime ye\cel{1}\sou{virtuala}\cel{1}: ube li}
\l{\cel{3}interrenkontrus se li esus plulongigita). - DEFIS.}
\l{}
\l{\sou{realgaro}. (\sou{kemio)}. Reda arseno-sulfajo.\cel{2}\sou{Simbolo kemiala}:}
\l{\cel{3}AsS.}
\l{}
\l{\sou{rebela}. Qua refuzas obediar la autoritati legitima. - DEFIS.}
\l{}
\l{\sou{rebordo}. (\sou{tekn.)}\cel{2}Bordo igita plu alta kam la centro-parto.}
\l{\cel{3}- FIRS.}
\l{}
\l{\sou{rebuso}. Esprito-ludo qua konsistas ye reprezentar vorto, fra-}
\l{\cel{3}zo per figuri di objekti di qui la nomi memorigas, per}
\l{\cel{3}homonimeso, ta vorto, ta frazo. - DEFIR.}
\l{}
\l{\sou{recensar}. (\sou{trans.)}\cel{2}(\sou{pri libri, artikli di revuo, e c.)}.}
\l{\cel{3}Analizar (por\cel{1}\sou{publico}) olia kontenajo, opinionar pri lia}
\l{\cel{3}valoro spiritala, e c. - DER.}
\l{}
\l{\sou{recenta}. Qua eventis, existeskis,nur poka tempo ante la ins-}
\l{\cel{3}tanto konsiderata. - EFIS.}
\l{}
\l{\sou{recepcionar}. (\sou{trans.}) Verifikar (pri qualeso e quanto) livrajo}
\l{\cel{3}por saveskar ka lu esas absolute konforma a la klauzi di}
\l{\cel{3}la kompro-kontrato. - F.}
\l{}
\l{\sou{recepto}. I. Indiko quan onu recevas pri la preparo di ula}
\l{\cel{3}medikamenti, di ula dishi. - II. (\sou{medicino}).Preskripto}
\l{\cel{3}da mediko per qua lu donas la formulo di medikamento. -}
\l{\cel{3}DEFIRS.}
\l{}
\l{\sou{receptaklo}. (\sou{bot}.) Parto di la floro ube su insertas la ka-}
\l{\cel{3}lico, la korolo, la stamini e, normale, la karpelo (o}
\l{\cel{3}karpeli) di la ovario. - EFS.}
\l{}
\l{\sou{recesiva.}\cel{2}(\sou{biol}.) Dicesas, pri la Mendel-lego (+\sou{leyo)}, por}
\l{\cel{3}qualifikar la karakteri qui dominacesas - t.e. qui trans-}
\l{\cel{3}misesas, a la hibrido od a la mestico - da un ek la pa-}
\l{\cel{3}tri, ma sen aparar en la genituro. - F.}
\l{}
\l{recevar. (\sou{trans., de)}. Esar la sendario di sendajo, la donario}
\l{\cel{3}di donajo, la donacario di donaco, e c.- esar la objekto}
\l{\cel{3}di ago. - EFIS.}
}
